\section{审判监督程序（再审）}
\begin{enumerate}
  \item 已经生效的裁判，无论是否执行。没有次数限制。
  \item \zd{提起}。
        \begin{enumerate}
          \item 申诉。不触发再审程序。
                \begin{enumerate}
                  \item 当事人及其法定代理人、近亲属、利害关系案外人、委托律师等。
                  \item 刑罚执行完毕 2 年内提出的申诉都应当受理；超过 2 年：可能无罪的、重大疑难还是可以受理的。
                  \item 找法院：
                        \begin{enumerate}
                          \item 原则\zd{终审法院}；越级：告知、交给、自己来；越多级没用；
                          \item 死刑找核准法院 \yjt 下级法院审查报告 \yjt 层报核准法院。
                          \item 2 级申诉制。死缠烂打，应当驳回 \hz 通知不予重新审判。
                        \end{enumerate}
                  \item 找检察院。终审法院同级检察院；越级：交给、自己来；也是 2 级申诉制，不再复查。
                  \item 无顺序要求。
                \end{enumerate}
          \item 法院。本院，提交审委会\zd{决定再审}；上级法院\zd{指令再审}（除非原审更适合）\hz \zd{决定提审}（法律错大疑复）。
          \item 检察院。上抗下。\zd{提出抗诉}，向谁抗谁来审，事实不清可指令。
        \end{enumerate}
  \item 重新审判。
        \begin{enumerate}
          \item 应当另行组成合议庭，原审人员回避。
          \item 级别。\zd{原来一审按一审，原来二审按二审，凡是提审统统按二审}。
          \item 应当制作再审决定书。
          \item 审理方式。
                \begin{enumerate}
                  \item 应当开庭。\zd{一抗事证加重罚}。
                        \begin{enumerate}
                          \item 依\zd{一}审程序审理；\zd{抗}诉。
                          \item 依二审程序需要对\zd{事}实和\zd{证}据进行审理的。
                          \item 可能对原审被告\zd{加重}刑\zd{罚}的。
                        \end{enumerate}
                  \item 可以不开庭。
                \end{enumerate}
          \item 撤回抗诉。
                \begin{enumerate}
                  \item 开庭审理前撤回，应当裁定准许。
                  \item 没有派员出庭 + 未说明理由，按撤回抗诉处理。
                \end{enumerate}
          \item 撤回申诉。
                \begin{enumerate}
                  \item 在再审期间撤回申诉，应当裁定准许；原判有错误的，应当不予准许。
                  \item 不到庭、退庭，应当裁定按撤回申诉处理；
                \end{enumerate}
          \item 重点审查原则。针对申诉、抗诉和决定再审的理由进行审理。必要时对事实、证据和适用法律全面审查。
        \end{enumerate}
  \item 原则不停止原判决、裁定的执行。但可能改判无罪 \hz 减轻而导致刑期届满，\zd{可以中止}执行。
  \item 强制措施谁启动谁决定。
  \item 下落不明 \hz 未到庭 \yjt 中止；到案后恢复；\zd{超过 2 年仍查无下落} \yjt 终止。
  \item 重新审理后的处理。
        \begin{enumerate}
          \item \zd{应当维持}。原判认定事实、适用法律正确、量刑适当。
          \item \zd{纠正维持}。认定事实、适用法律\zd{有瑕疵}。
          \item \zd{应当改判}。法律错 \hz 量刑错。
          \item 原判事实不清 \hz 证据不足 \yi 按二审审可改可发；\er 按一审，该咋判咋判。
          \item 审限 3 + 3
        \end{enumerate}
\end{enumerate}
